% =============================================================================
\section{基础：vLLM Semantic Router 研究计划}
\label{sec:foundation}
% =============================================================================

\Cref{tab:ownwork} 列出了支撑本文的那些论文。我们将它们分成四个支柱。

\begin{table}[t]
\centering
\caption{构成 WRP 愿景基础的 vLLM Semantic Router 项目论文。每一行都标明该论文主要涉及哪些 WRP 维度。 \textbf{W}=Workload，\textbf{R}=Router，\textbf{P}=Pool。}
\label{tab:ownwork}
\small
\renewcommand{\arraystretch}{1.15}
\begin{tabularx}{\textwidth}{l X c c c}
\toprule
\textbf{简称} & \textbf{标题 / 发表渠道} & \textbf{W} & \textbf{R} & \textbf{P} \\
\midrule
\rowcolor{lightgray}
PoolRouting~\cite{poolrouting2026}
  & Token-Budget-Aware Pool Routing（arXiv） & $\bullet$ & $\bullet$ & $\bullet$ \\
FleetOpt~\cite{fleetopt2026}
  & Analytical Fleet Provisioning with C\&R（arXiv） & $\bullet$ & $\bullet$ & $\bullet$ \\
\rowcolor{lightgray}
1/W Law~\cite{onewlaw2026}
  & Context-Length Routing and Energy Efficiency（arXiv） & $\bullet$ & $\bullet$ & $\bullet$ \\
\midrule
FeedbackDet~\cite{feedbackdet2026}
  & mmBERT-32K Feedback Detector（HuggingFace） & $\bullet$ & $\bullet$ &  \\
\rowcolor{lightgray}
AVR~\cite{avr2026}
  & Adaptive VLM Routing for CUAs（arXiv） & $\bullet$ & $\bullet$ &  \\
MPExt~\cite{mpext2025}
  & Multi-Provider Extensions for Agentic AI（IETF） & $\bullet$ & $\bullet$ &  \\
\rowcolor{lightgray}
OpenClaw~\cite{openclaw2026}
  & Personal assistant + Gateway（GitHub OSS） & $\bullet$ & $\bullet$ &  \\
SemCache~\cite{wang2025cache}
  & Category-Aware Semantic Caching（arXiv） &  & $\bullet$ & $\bullet$ \\
FleetSim~\cite{fleetsim2026}
  & Queueing-Grounded Fleet Capacity Planner（arXiv） & $\bullet$ &  & $\bullet$ \\
\midrule
\rowcolor{lightgray}
\vsr{}~\cite{vllmsr2026}
  & Signal-Driven Decision Routing for MoM Models（arXiv） &  & $\bullet$ &  \\
ProbPol~\cite{probpol2026}
  & Conflict-Free Policy for Probabilistic ML Predicates（arXiv） &  & $\bullet$ &  \\
\rowcolor{lightgray}
FastRouter~\cite{fastrouter2026}
  & 98$\times$ Faster LLM Routing（arXiv） &  & $\bullet$ &  \\
WhenToReason~\cite{wang2025reason}
  & Semantic Router for vLLM（NeurIPS MLForSys） &  & $\bullet$ &  \\
\rowcolor{lightgray}
mmBERT-Embed~\cite{mmberted2025}
  & 2D Matryoshka Embeddings for Routing（HuggingFace） &  & $\bullet$ &  \\
HaluGate~\cite{halugate2025}
  & Token-Level Hallucination Detection（vLLM Blog） &  & $\bullet$ &  \\
\rowcolor{lightgray}
FactCheck~\cite{factcheck2026}
  & mmBERT-32K Factcheck Classifier（HuggingFace） &  & $\bullet$ &  \\
OATS~\cite{oats2026}
  & Outcome-Aware Tool Selection（arXiv） &  & $\bullet$ &  \\
\rowcolor{lightgray}
VCD~\cite{vcd2026}
  & Visual Confused Deputy: CUA Security（arXiv） &  & $\bullet$ &  \\
SafetyL1~\cite{safetyl1_2026}
  & MLCommons Binary Safety Classifier（HuggingFace） &  & $\bullet$ &  \\
\rowcolor{lightgray}
SafetyL2~\cite{safetyl2_2026}
  & MLCommons 9-Class Hazard Classifier（HuggingFace） &  & $\bullet$ &  \\
SIRP~\cite{sirp2025}
  & Semantic Inference Routing Protocol（IETF） &  & $\bullet$ &  \\
\bottomrule
\end{tabularx}
\end{table}

\subsection{支柱 1：路由架构}
\label{sec:foundation:routing}

\vsr{} 立场论文~\cite{vllmsr2026}提出了由信号驱动的决策路由：把异构信号，例如关键词模式、嵌入相似度、领域分类器和语言检测，组合成部署特定的路由策略，总共覆盖十三类信号。ProbPol~\cite{probpol2026}形式化了当概率式机器学习信号同时触发时的冲突检测问题；mmBERT-Embed~\cite{mmberted2025}则提供了嵌入骨干，并通过二维 Matryoshka 结构支持可配置的延迟-质量权衡（见 \Cref{sec:router:static}）。

一个设计约束是，路由器必须同时满足\emph{低延迟、低资源占用}：它不应要求专用 GPU，因为那个 GPU 本可以直接拿来承载 LLM 推理。FastRouter~\cite{fastrouter2026}正是围绕这一点展开：面对长上下文请求，它把延迟降低了 98$\times$，同时 GPU 占用低于 800\,MB（见 \Cref{sec:router:static}）。

除请求发出时的分发以外，路由栈还包含响应时机制。WhenToReason~\cite{wang2025reason}负责判断某个查询是否需要进入 reasoning 模式。SemCache~\cite{wang2025cache}通过带有领域特定相似度阈值的类别感知语义缓存，消除重复推理（见 \Cref{sec:workload:warm-cold}）。Feedback Detector~\cite{feedbackdet2026}使得系统可以根据每轮用户满意度信号进行在线路由自适应（见 \Cref{sec:router:online}）。HaluGate~\cite{halugate2025}与 Factcheck Classifier~\cite{factcheck2026}进一步把路由扩展到响应阶段验证：它们把事实性查询导向更可靠的模型，并把每个模型的幻觉率反馈回路由策略中（见 \Cref{sec:router:static}）。

该路由栈还通过遵循 MLCommons AI Safety Taxonomy 的分层分类流水线来执行内容安全与隐私控制。Level-1 二元分类器~\cite{safetyl1_2026}对每个入站请求做一次安全/不安全判断；被标记为不安全的请求会升级到 Level-2 九类危害分类器~\cite{safetyl2_2026}，后者将其映射到具体类别（暴力犯罪、隐私、错误信息等）。这一两阶段设计让良性流量几乎不增加额外开销，同时仍能支持按类别执行策略，例如直接阻断侵犯隐私的提示、将易产生错误信息的查询导向事实性更强的模型、或者对专业建议类请求进行合规审计记录（见 \Cref{sec:router:static}）。

\subsection{支柱 2：资源池路由与集群优化}
\label{sec:foundation:fleet}

令牌预算资源池路由~\cite{poolrouting2026}是起点：它根据估计的总令牌预算，将请求分发到短上下文池或长上下文池，在生产轨迹上减少了 17--39\% 的 GPU 实例（见 \Cref{sec:router:static}）。

FleetOpt~\cite{fleetopt2026}把资源池路由与网关层压缩统一到同一个分析框架中。给定工作负载 CDF 与延迟目标，它可以推导出最小成本的双池集群，以及最优边界 $\boundary^*$ 与最优压缩参数 $\gamma^*$。与之互补，inference-fleet-sim~\cite{fleetsim2026}使用覆盖 A10G、A100 与 H100 的离散事件仿真器验证了这些结果（见 \Cref{sec:pool:homo-hetero}）。

1/W 定律~\cite{onewlaw2026}则量化了能耗维度：路由拓扑比 GPU 代际更能撬动能效，而且这种收益是可乘的（见 \Cref{sec:pool:energy}）。

\subsection{支柱 3：多模态与智能体路由}
\label{sec:foundation:agent}

Adaptive VLM Routing（AVR）~\cite{avr2026}把路由扩展到视觉语言 CUA 任务，提出了一个三层框架：第 1 层是多模态分类器（预路由约 $\sim$10\,ms），第 2 层是带 token 过滤评分的小型 VLM 置信度探测，第 3 层是大模型升级。AVR 预计可将成本降低多达 78\%，同时与全程使用大模型的基线相比仅落后 2 个百分点以内。

OATS~\cite{oats2026}处理的是关键请求路径中的工具选择问题：通过离线方式把工具嵌入插值到“成功查询”质心方向，它在不增加服务时开销的情况下，将 MetaTool 上的 NDCG@5 从 0.869 提高到 0.940。

Visual Confused Deputy~\cite{vcd2026}则把 CUA 感知错误重构为一种安全漏洞，并提出双通道对比分类方法，分别评估点击目标与推理文本，达到 F1$=$0.915。与 AVR 结合后，VCD 使得同一套路由流水线同时提供效率与安全。

\subsection{支柱 4：治理与标准}
\label{sec:foundation:governance}

在标准层面，Semantic Inference Routing Protocol（SIRP）~\cite{sirp2025}在 IETF 中为 AI 推理系统中的内容级分类与语义路由规定了一套框架。Multi-Provider Extensions~\cite{mpext2025}则把智能体 AI 推理 API 扩展到了多提供方环境。
